% =====================================================================
%  U14：PDF p.315--p.325（原书印刷 p.305--p.315，页码以页面页脚为准）
%  第 10 章章首 + §10.1 Opers 与局部系统 + §10.2 Harish--Chandra 范畴
%  （10.2.1--10.2.3）
%  注：任务书的页段标为印刷 301--311，但页面页脚显示 PDF 315 = 印刷 305
%  （本章起偏移为 PDF − 10）；批注中的 p.N 一律用页脚所示的印刷页码。
% =====================================================================

\chapter{构造 Langlands 对应}

% [Tutor Note: p.305。本章把第 1 章的问题与第 2--9 章关于临界水平中心的结果接起来：参数一侧是 $\Loc_{{}^L G}(D^\times)$，被参数化一侧是 $\hat{\fg}_\kappa\text{-}\mathrm{mod}$；关键手段是把范畴沿中心特征的谱“纤维化”。首式 (10.0-25) 就是第 1 章 (1.2-7) 的复述，但原书在此另给了一个编号 (10.0-25)：第 10 章引言位于 §10.1 之前，故编号中节号为 0，请读者不要去第 1 章找 (10.0-25)。]
该是用前几章所得结果来推进圈群的局部 Langlands 对应的时候了。回到第 1 章末尾提出的问题：令
\begin{equation}
    \Loc_{{}^L G}(D^\times)
    =\left\{\partial_t+A(t),\ A(t)\in{}^L\fg((t))\right\}\Big/{}^L G((t))
    \tag{10.0-25}
\end{equation}
是去心圆盘 $D^\times=\Spec\C((t))$ 上 ${}^L G$-联络的规范等价类集合。我们在第 1 章已论证过，应当把 $\Loc_{{}^L G}(D^\times)$ 当作圈群 $G((t))$ 的 Langlands 参数空间。回忆圈群 $G((t))$ 作用在水平 $\kappa$ 的（光滑）$\hat{\fg}$-模范畴 $\hat{\fg}_\kappa\text{-}\mathrm{mod}$ 上（该范畴的定义见 1.3.6 节）。我们提出了如下问题：

\begin{quote}
\itshape
对每个局部 Langlands 参数 $\sigma\in\Loc_{{}^L G}(D^\times)$，配上一个子范畴 $\hat{\fg}_\kappa\text{-}\mathrm{mod}_\sigma\subset\hat{\fg}_\kappa\text{-}\mathrm{mod}$，它在圈群 $G((t))$ 的作用下稳定。
\end{quote}

更有雄心地说，我们希望把范畴 $\hat{\fg}_\kappa\text{-}\mathrm{mod}$ 表示成“纤维化”于局部 Langlands 参数空间 $\Loc_{{}^L G}(D^\times)$ 之上，各范畴 $\hat{\fg}_\kappa\text{-}\mathrm{mod}_\sigma$ 是“纤维”，而群 $G((t))$ 沿这些纤维作用。若能做到这一点，我们就把这个纤维化看作范畴 $\hat{\fg}_\kappa\text{-}\mathrm{mod}$ 在 $\Loc_{{}^L G}(D^\times)$ 上的“谱分解”。

在第 2 章开头，我们提出了解决这个问题的一种可能方案。我们指出，任何阿贝尔范畴都可以看成“纤维化”于其中心的谱之上。于是我们的想法是：描述范畴 $\hat{\fg}_\kappa\text{-}\mathrm{mod}$ 的中心（对每个 $\kappa$ 值），并考察它的谱是否与 Langlands 参数空间 $\Loc_{{}^L G}(D^\times)$ 有关。

% [Tutor Note: p.306。这一步“中心”是范畴论意义上的中心（与所有函子自然交换的自同态环），它被等同于完备泛包络代数中心的像。关键事实：$\kappa\neq\kappa_c$ 时 $Z_\kappa(\hat{\fg})=\C$（命题 4.3.9），中心是平凡的，谱只有一点，方案失效；$\kappa=\kappa_c$ 时中心由 Feigin--Frenkel 的结果可知非常大，才有非平凡的谱可纤维化。这个对比是整章的出发点。]
我们已把范畴 $\hat{\fg}_\kappa\text{-}\mathrm{mod}$ 的中心等同于结合代数 $\widetilde{U}_\kappa(\hat{\fg})$ 的中心 $Z_\kappa(\hat{\fg})$，而 $\widetilde{U}_\kappa(\hat{\fg})$ 是水平 $\kappa$ 的 $\hat{\fg}$ 完备泛包络代数，定义见 2.1.2 节。接着我们描述了代数 $Z_\kappa(\hat{\fg})$。根据命题 4.3.9，若 $\kappa\neq\kappa_c$（临界水平），则 $Z_\kappa(\hat{\fg})=\C$。因此我们的方案对 $\kappa\neq\kappa_c$ 行不通。但我们发现，临界水平的中心 $Z_{\kappa_c}(\hat{\fg})$ 非常不平凡，且确实与去心圆盘上的 ${}^L G$-联络有关。

在本章中，我们将依照与 D. Gaitsgory 的合作工作 [Frenkel and Gaitsgory 2004; Frenkel and Gaitsgory 2006c; Frenkel and Gaitsgory 2006a; Frenkel and Gaitsgory 2005; Frenkel and Gaitsgory 2006b; Frenkel and Gaitsgory 2007c]，利用这些结果对圈群的局部 Langlands 对应提出更精确的猜想，并为这些猜想提供一些证据。我们还将讨论这些猜想对整体几何 Langlands 对应的影响。

% [Tutor Note: p.306。作者在此提到 A. Beilinson 的另一套方案 [2006]：用低于临界水平的仿射 Kac--Moody 代数表示构造局部几何 Langlands 对应。两套方案的关系（“低于临界”与“临界”、表示上的差别）至今没有被完全理清，这一点值得留意，因为临界水平在本章扮演不可替代的角色。]
我们要指出，A. Beilinson 对局部几何 Langlands 对应另有一套方案 [2006]，用的是水平小于临界水平的仿射 Kac--Moody 代数的表示。弄清他的方案与我们的方案之间的联系，会很有意思。

下面是本章更详细的计划。在 10.1 节，我们回顾局部系统与 opers 的关系。我们期待圈群真正的 Langlands 参数空间是去心圆盘 $D^\times$ 上 ${}^L G$-局部系统的空间。另一方面，$\hat{\fg}_{\kappa_c}$-模范畴纤维化于 $D^\times$ 上 ${}^L G$-opers 的空间之上。我们将看到这带来一些不平凡的推论。接着，在 10.2 节我们引入 Harish--Chandra 范畴；它们是紧子群 $K\subset G(F)$ 下不变向量子空间的范畴类比物。我们的纲领是用 Langlands 对偶群 ${}^L G$ 来描述 $\hat{\fg}_{\kappa_c}$ 上的 Harish--Chandra 模范畴，并在经典 Langlands 对应中的相应对象之间找出联系与平行。我们在两个重要例子中实施这一纲领：“无分歧”情形，即 $D^\times$ 上的局部系统同构于平凡局部系统；以及“驯顺分歧”情形，即相应的联络有正则奇点与幂幺单值性。

% [Tutor Note: p.306--307。这段是本章的路线图，也是全书其余部分的纲领：10.3 无分歧（Satake 对应的范畴化）；10.4 驯顺分歧（仿射 Hecke 代数、Springer 纤维）；10.5 从局部到整体（Hecke 本征层）。阅读时建议先抓住两个对照：局部系统 ↔ opers（参数一侧的两种版本），$\hat{\fg}_{\kappa_c}$-模范畴 ↔ Harish--Chandra 范畴（被参数化一侧的两层）。]
无分歧情形在 10.3 节详细讨论。我们先回忆，在经典情形它由 Satake 对应表达：$G(F)$ 的无分歧表示的同构类与 Langlands 对偶群 ${}^L G$ 中的半单共轭类相对应。然后我们讨论这一对应的范畴版本。对与平凡局部系统相联系的圈群范畴表示，我们给出两种实现：一种用 $\hat{\fg}_{\kappa_c}$-模，另一种用仿射 Grassmann 流形上的 D-模；并给出它们之间的猜想等价。特别地，我们要证明相应的 Harish--Chandra 模范畴等价于向量空间范畴，这与经典图像一致。

驯顺分歧情形在 10.4 节讨论。在经典情形，$G(F)$ 的相应不可约表示与 $G(F)$ 的仿射 Hecke 代数的不可约模一一对应。根据 [Kazhdan and Lusztig 1987; Chriss and Ginzburg 1997]，这些模可以实现在 ${}^L G$ 的幂幺元所对应的 Springer 纤维的代数 $K$-理论中作为子商（这些幂幺元在几何情形中扮演局部系统单值性的角色）。这提示：在范畴版本中，$\hat{\fg}_{\kappa_c}$ 上相应的 Harish--Chandra 模范畴应当等价于 Springer 纤维上适当的凝聚层范畴（更精确地说，这里应当考虑导出范畴）。这样一个确切猜想已在 [Frenkel and Gaitsgory 2006c] 中给出。在 10.4 节，我们将依照 [Frenkel and Gaitsgory 2006c; Frenkel and Gaitsgory 2006a; Frenkel and Gaitsgory 2005; Frenkel and Gaitsgory 2006b] 讨论这个猜想及各种支持性证据。我们指出，Springer 纤维的出现可以追溯到上一章所详述的“把 Wakimoto 模配给 Miura oper”的构造。我们强调经典图像与几何图像之间的平行。我们还给出 $\fg=\fsl_2$ 情形的一些显式计算，用来说明一般理论。

% [Tutor Note: p.307。注意这里两处人名与术语按术语表保留英文：affine Hecke algebra、Springer fiber、Miura oper、Wakimoto module。驯顺分歧情形在几何上对应“正则奇点 + 幂幺单值性”，幂幺轨道经 Springer 对应联系到 Springer 纤维，这是 L-包分类中的标准桥（参看 Lusztig 的 Springer 表示理论）。]
最后，在 10.5 节我们说明这些结果对整体几何 Langlands 对应意味着什么。我们展示局部化函子如何把 $\hat{\fg}_{\kappa_c}$ 上局部 Harish--Chandra 模范畴与带抛物结构的曲线上 $G$-丛模叠的整体 Hecke 本征层范畴联系起来。我们期待，至少在一般情形，局部范畴与整体范畴彼此等价。因此，我们关于局部范畴的结果与猜想就翻译成关于 Hecke 本征层范畴结构的明确陈述。这样，仿射 Kac--Moody 代数的表示论就为整体 Langlands 对应提供了有价值的洞见。\transnote{这里的“局部化函子”指把 $\operatorname{Bun}_G$ 上的 D-模与临界水平的 $\hat\fg_{\kappa_c}$-模联系起来的函子（[Beilinson and Drinfeld 1997] 的局部化，以及 [Frenkel and Gaitsgory 2005] 在仿射 Grassmann 流形上的版本）。整体几何 Langlands 的构造方向是由局部系统（即 Hecke 本征值）得到 Hecke 本征层，而不是反过来；本节只给纲领，细节见 10.5 节。[来源：原书 p.307；[Beilinson and Drinfeld 1997]；[Frenkel and Gaitsgory 2005]。]}

\begin{keypoints}
    \item \textbf{核心思想：} 用临界水平的范畴 $\hat{\fg}_{\kappa_c}\text{-}\mathrm{mod}$ 沿 opers 空间的纤维化，来实现圈群的局部 Langlands 对应；真正的参数空间是局部系统空间 $\Loc_{{}^L G}(D^\times)$，而范畴自然纤维化于其“opers 版本” $\Op_{{}^L G}(D^\times)$ 之上。
    \item \textbf{数学工具：} 临界水平中心 $Z_{\kappa_c}(\hat{\fg})\cong\operatorname{Fun}\Op_{{}^L G}(D^\times)$；沿概形的范畴基变换与笛卡尔方块；Harish--Chandra 范畴与范畴 Hecke 代数。
    \item \textbf{关键结论：} 遗忘映射 $\alpha:\Op_{{}^L G}(D^\times)\to\Loc_{{}^L G}(D^\times)$ 猜测为满射（猜想 10.1.1）；$\hat{\fg}_{\kappa_c}\text{-}\mathrm{mod}$ 应是猜想中万有范畴在 $\Op$ 上的基变换（(10.1-3)）；同一参数纤维内的不同 opers 给出相互等价的范畴（猜想 10.1.2）。
    \item \textbf{记号提醒：} $\kappa_c$ 为临界水平；${}^L B\subset{}^L G$ 是固定的 Borel 子群，${}^L\mathrm{b}$ 是其李代数；$R^{\pi(K)}$ 表示 $K$-不变向量空间；$\hat{\fg}_{\kappa_c}\text{-}\mathrm{mod}_\chi$ 表示固定中心特征 $\chi$ 的纤维范畴。
\end{keypoints}

\section{Opers 与局部系统}

\begin{keypoints}
    \item \textbf{核心思想：} 真正的 Langlands 参数是 $D^\times$ 上的局部系统（二元组），而临界水平的范畴只自然纤维化于 opers 空间（三元组）之上；连接两者的遗忘映射 $\alpha$ 是整章的枢纽。
    \item \textbf{数学工具：} 中心同构 $Z_{\kappa_c}(\hat{\fg})\cong\operatorname{Fun}\Op_{{}^L G}(D^\times)$ 与中心特征标；opers 的规范类描述 (10.1-1)；范畴沿概形的基变换与笛卡尔方块 (10.1-3)。
    \item \textbf{关键结论：} $\alpha$ 猜测为满射（Deligne 已证 $SL_n$ 情形）；$\hat{\fg}_{\kappa_c}\text{-}\mathrm{mod}$ 应是万有范畴 $\cC$ 沿 $\alpha$ 的基变换；同一局部系统纤维内不同 opers 给出的范畴相互等价（猜想 10.1.1、10.1.2）。
    \item \textbf{记号提醒：} 局部系统 $(\cF,\nabla)$ 对比 oper $(\cF,\nabla,\cF_{{}^L B})$；$\ell$ 是 ${}^L\fg$ 的秩（单根个数）；${}^L\mathrm{b}$ 是 ${}^L B$ 的李代数。
\end{keypoints}

% [Tutor Note: p.307。本节的核心是“两个参数空间”的对照：$\Loc$（局部系统，二元组）与 $\Op$（opers，三元组），以及遗忘映射 $\alpha$。整个第 10 章的对应都建立在 $\alpha$ 的满射性（猜想 10.1.1）与“范畴沿 $\alpha$ 基变换”（(10.1-3)）之上。]
在好几章的篇幅里，我们详细描述了临界水平的中心 $Z_{\kappa_c}(\hat{\fg})$。根据定理 4.3.6 以及第 8 章得到的相关结果，$Z_{\kappa_c}(\hat{\fg})$ 同构于去心圆盘 $D^\times$ 上 ${}^L G$-opers 空间上的函数代数 $\operatorname{Fun}\Op_{{}^L G}(D^\times)$，并且这个同构与两个代数上的各种对称性与结构相容。现在注意，点 $\chi\in\Op_{{}^L G}(D^\times)$ 与同态（等价地，特征标）
\[
    \operatorname{Fun}\Op_{{}^L G}(D^\times)\longrightarrow\C
\]
一一对应，后者就是“在 $\chi$ 处取值”这个求值映射。因此，$\Op_{{}^L G}(D^\times)$ 的点参数化了中心特征标 $Z_{\kappa_c}(\hat{\fg})\to\C$。\transnote{“与各种对称性相容”的确切形式见第 8 章：该同构把 $\operatorname{Fun}\Op_{{}^L G}(D^\times)$ 与中心 $Z_{\kappa_c}(\hat\fg)$ 的结构对应起来，从而把中心的特征标等同于 opers 上的点，这是本节全部构造的出发点。同构的源头是 [Feigin and Frenkel 1992] 对临界水平中心（Gelfand--Dikiĭ 代数）的计算。[来源：原书 p.307；第 8 章；[Feigin and Frenkel 1992]。]}

给定一个 ${}^L G$-oper $\chi\in\Op_{{}^L G}(D^\times)$，我们定义范畴
\[
    \hat{\fg}_{\kappa_c}\text{-}\mathrm{mod}_\chi
\]
为 $\hat{\fg}_{\kappa_c}\text{-}\mathrm{mod}$ 的满子范畴，其对象是临界水平的 $\hat{\fg}$-模（从而是 $\widetilde{U}_{\kappa_c}(\hat{\fg})$-模），并且中心 $Z_{\kappa_c}(\hat{\fg})\subset\widetilde{U}_{\kappa_c}(\hat{\fg})$ 按对应于 $\chi$ 的中心特征标作用。由于代数 $\Op_{{}^L G}(D^\times)$ 作用在范畴 $\hat{\fg}_{\kappa_c}\text{-}\mathrm{mod}$ 上，我们可以说范畴 $\hat{\fg}_\kappa\text{-}\mathrm{mod}$“纤维化”于空间 $\Op_{{}^L G}(D^\times)$ 之上，使得对应于 $\chi\in\Op_{{}^L G}(D^\times)$ 的纤维范畴是 $\hat{\fg}_{\kappa_c}\text{-}\mathrm{mod}_\chi$。\transnote{此处原书 p.308 把范畴印作 $\hat\fg_\kappa\text{-}\mathrm{mod}$（下标无 $c$），而同句的纤维范畴、以及下一段“acts on \dots\ the category”都印作 $\hat\fg_{\kappa_c}$；按上下文应读作临界水平 $\kappa_c$。译文照原文保留了这个 $\kappa$，读时不必追究。[来源：原书 p.308（PDF p.318）页图。]}

% [Tutor Note: p.308 易错处。“中心按中心特征标作用”给出的是 $\hat{\fg}_{\kappa_c}\text{-}\mathrm{mod}$ 按 $\Spec Z_{\kappa_c}(\hat{\fg})$ 的直和分解（广义特征空间分解），这只在中心作用半单时才是直和；一般情形得到的是“纤维化”意义上的分解。作者用“fibering”而不用“decomposition”，正是为了避免声称直和分解。另注意原文此处写的是 $\hat{\fg}_\kappa\text{-}\mathrm{mod}$（用 $\kappa$ 而非 $\kappa_c$），按上下文应理解为临界水平，译文照原文保留。]
更一般地，对任意闭代数子簇 $Y\subset\Op_{{}^L G}(D^\times)$（不必是单点），有理想
\[
    I_Y\subset\operatorname{Fun}\Op_{{}^L G}(D^\times)\simeq Z_{\kappa_c}(\hat{\fg}),
\]
由那些在 $Y$ 上取零的函数组成。于是有满子范畴 $\hat{\fg}_{\kappa_c}\text{-}\mathrm{mod}_Y\subset\hat{\fg}_{\kappa_c}\text{-}\mathrm{mod}$，其对象是临界水平的 $\hat{\fg}$-模，且 $I_Y$ 作用为零。这个范畴是范畴 $\hat{\fg}_{\kappa_c}\text{-}\mathrm{mod}$ 沿态射 $Y\to\Op_{{}^L G}(D^\times)$ 作“基变换”的一个例子。容易把这个定义推广到任意仿射概形 $Y$ 连同态射 $Y\to\Op_{{}^L G}(D^\times)$。相应的基变换范畴 $\hat{\fg}_{\kappa_c}\text{-}\mathrm{mod}_Y$ 随后可以“粘合”起来，这使我们对任意映入 $\Op_{{}^L G}(D^\times)$ 的概形 $Y$ 定义基变换范畴 $\hat{\fg}_{\kappa_c}\text{-}\mathrm{mod}_Y$。不难把这推广为“阿贝尔范畴纤维化于代数叠之上”的一般概念（见 [Gaitsgory 2005]）。范畴纤维化于代数叠 $\cY$ 之上的一个模型例子，是代数叠 $\cX$ 上凝聚层的范畴，其中 $\cX$ 带有态射 $\cX\to\cY$。

% [Tutor Note: p.308 含糊处。“基变换范畴”在这里不是通常的纤维积，而是“模掉 $I_Y$ 作用”得到的满子范畴；作者随后说可以粘合、推广到任意概形与叠，这需要 [Gaitsgory 2005] 中的形式化（范畴化版本的“拟凝聚层沿态射拉回”）。读者若只关心主定理，可把 $Y$ 取为单点 $\chi$，此时 $I_Y$ 就是极大理想。]
回忆群 $G((t))$ 作用在 $\widetilde{U}_{\kappa_c}(\hat{\fg})$ 上，也作用在范畴 $\hat{\fg}_{\kappa_c}\text{-}\mathrm{mod}$ 上。根据命题 4.3.8，$G((t))$ 在 $Z_{\kappa_c}(\hat{\fg})\subset\widetilde{U}_{\kappa_c}(\hat{\fg})$ 上的作用是平凡的。因此子范畴 $\hat{\fg}_{\kappa_c}\text{-}\mathrm{mod}_\chi$（更一般地，$\hat{\fg}_{\kappa_c}\text{-}\mathrm{mod}_Y$）在 $G((t))$ 作用下稳定。于是群 $G((t))$ 沿着“纤维化” $\hat{\fg}_{\kappa_c}\text{-}\mathrm{mod}\to\Op_{{}^L G}(D^\times)$ 的纤维作用（详见 [Frenkel and Gaitsgory 2006c] 第 20 节）。

纤维化 $\hat{\fg}_{\kappa_c}\text{-}\mathrm{mod}\to\Op_{{}^L G}(D^\times)$ 几乎给出了我们想要的圈群局部 Langlands 对应。但有一个重要差别：我们要求范畴 $\hat{\fg}_{\kappa_c}\text{-}\mathrm{mod}$ 纤维化于 $D^\times$ 上局部系统的空间 $\Loc_{{}^L G}(D^\times)$ 之上；然而我们已证明，$\hat{\fg}_{\kappa_c}\text{-}\mathrm{mod}$ 纤维化于 ${}^L G$-opers 的空间 $\Op_{{}^L G}(D^\times)$ 之上。

这两个空间的区别何在？${}^L G$-局部系统是一个二元组 $(\cF,\nabla)$，其中 $\cF$ 是 ${}^L G$-丛、$\nabla$ 是 $\cF$ 上的联络；而 ${}^L G$-oper 是一个三元组 $(\cF,\nabla,\cF_{{}^L B})$，其中 $\cF$ 与 $\nabla$ 如前，$\cF_{{}^L B}$ 是附加的结构：把 $\cF$ 约化到（固定的）Borel 子群 ${}^L B\subset{}^L G$，并满足 4.2.1 节所解释的横截性条件。因此对任意曲线 $X$，显然有遗忘映射
\[
    \Op_{{}^L G}(X)\longrightarrow\Loc_{{}^L G}(X).
\]
该映射在 $(\cF,\nabla)\in\Loc_{{}^L G}(X)$ 上的纤维，由 $\cF$ 的所有满足关于 $\nabla$ 的横截性条件的 ${}^L B$-约化组成。

% [Tutor Note: p.309。通篇最容易混淆的一点：oper 的附加结构是“约化 $\cF_{{}^L B}$”，而不是联络的额外性质；遗忘映射忘掉的正是这个约化。横截性条件（4.2.1 节）保证局部可写成 (10.1-1) 的形状，从而 $\psi_i(t)$ 非零是“oper”区别于普通联络的本质条件。]
这个映射很可能不是满射，即它在某个特定局部系统 $(\cF,\nabla)$ 上的纤维可能是空的。例如，若 $X$ 是射影曲线且 ${}^L G$ 是伴随型群，则存在唯一的 ${}^L G$-丛 $\cF_0$，使得 $(\cF_0,\nabla)$ 上的纤维非空。此外，对这个 $\cF_0$，$(\cF_0,\nabla)$ 上的纤维对任意联络 $\nabla$ 都只含一个点（换言之，对 $\cF_0$ 上每个联络 $\nabla$，存在唯一的 ${}^L B$-约化 $\cF_{{}^L B}$ 满足关于 $\nabla$ 的横截性条件）。

当 $X=D^\times$ 时情形相当不同。此时任意 ${}^L G$-丛 $\cF$ 都可平凡化。因此联络 $\nabla$ 可以表示为一阶算子 $\partial_t+A(t)$，$A(t)\in{}^L\fg((t))$。但 $\cF$ 的平凡化不唯一；两个平凡化相差 ${}^L G((t))$ 中的一个元素。所以二元组 $(\cF,\nabla)$ 的等价类集合等同于商 (10.0-25)。

现在设 $(\cF,\nabla)$ 带有 oper 约化 $\cF_{{}^L B}$。那么只考虑那些来自 $\cF_{{}^L B}$ 的平凡化的 $\cF$ 的平凡化。这样的平凡化更少，因为两个平凡化现在相差 ${}^L B((t))$ 中的元素而非 ${}^L G((t))$ 中的元素。由于 oper 的横截性条件，联络 $\nabla$ 相对于任一这样的平凡化必有特殊形式：
\[
    \nabla=\partial_t+\sum_{i=1}^{\ell}\psi_i(t)f_i+\mathbf{v}(t),
\]
其中每个 $\psi_i(t)\neq0$，且 $\mathbf{v}(t)\in{}^L\mathrm{b}((t))$（见 4.2.2 节）。于是我们得到 opers 空间的一个具体实现，即作为规范等价类的空间：
\begin{equation}
\begin{aligned}
    \Op_{{}^L G}(D^\times)
    &=\Bigl\{\partial_t+\sum_{i=1}^{\ell}\psi_i(t)f_i+\mathbf{v}(t),\\
    &\qquad\qquad \psi_i\neq0,\ \mathbf{v}(t)\in{}^L\mathrm{b}((t))\Bigr\}\Big/{}^L B((t)).
\end{aligned}
    \tag{10.1-1}
\end{equation}

% [Tutor Note: p.309 易错处。(10.1-1) 与 4.2.2 节的 oper 定义必须对照阅读：$\ell$ 是单根个数（即 ${}^L\fg$ 的秩），$f_i$ 是 Chevalley 生成元中的根向量，$\psi_i(t)$ 是各简单根方向的系数，非零条件即横截性。注意 $\mathbf{v}(t)$ 取值于 Borel 李代数 ${}^L\mathrm{b}$，而不是整个 ${}^L\fg$；原书此处排的是正体 $\mathrm{b}$（不是通常的哥特体 $\mathfrak{b}$），译文照原书保留。若把 $\mathbf{v}$ 误取为 ${}^L\fg((t))$，条件就退化掉了。]
现在映射
\[
    \alpha:\Op_{{}^L G}(D^\times)\longrightarrow\Loc_{{}^L G}(D^\times)
\]
只是把形如 (10.1-1) 的算子的 ${}^L B((t))$-等价类，送到它的 ${}^L G((t))$-等价类。\transnote{上面的 (10.1-1) 正是 Drinfeld--Sokolov 意义下 oper 的规范类描述：$\ell$ 是单根个数，$f_i$ 是各简单根方向的 Chevalley 生成元，$\psi_i(t)\ne0$ 是 4.2.1 节横截性条件在此平凡化下的表现。对 ${}^L G=SL_2$，经 ${}^L B$-规范可把唯一的系数 $\psi_1$ 归一为 $1$，于是 oper 就是 $\partial_t+f_1+v(t)e_1$ 的形式，等价于去心圆盘上的一个射影结构；opers 与 KdV 方程族的联系正来自这种规范形。[来源：原书 p.309；[Drinfeld and Sokolov 1984]；[Beilinson and Drinfeld 2005]。]}

与上面讨论的射影曲线 $X$ 的情形不同，我们期待 $\alpha$ 对任意单李群 ${}^L G$ 都是满射。在 $SL_n$ 的情形这已由 [Deligne 1970] 证明，我们猜想一般情形也成立。\transnote{$\alpha$ 的纤维正是同一局部系统上的全部 oper 结构：给定 $(\cF,\nabla)$，纤维由满足横截性条件的 ${}^L B$-约化组成，故满射性等价于“$D^\times$ 上每个 ${}^L G$-局部系统都带有 oper 结构”。原书把 $SL_n$ 情形归功于 [Deligne 1970]，即 LNM 163《\'Equations diff\'erentielles \`a points singuliers r\'eguliers》（正则奇点线性方程）；对一般单群，这仍是猜想 10.1.1 的内容。[来源：原书 p.309--310；[Deligne 1970]。]}

\begin{conjecture}\label{conj:u14-alpha}
    映射 $\alpha$ 对任意单李群 ${}^L G$ 都是满射。
\end{conjecture}

% [Tutor Note: p.310 背景。$\alpha$ 的满射性等价于：$\Op$ 上的每个局部系统都能配上满足横截性的 Borel 约化。$SL_n$ 情形是 Deligne 的经典结果（与射影结构的 opers 描述有关）；一般单群情形远非显然，因为纤维可能为空（如正文所述伴随型群在射影曲线上的例子，但 $D^\times$ 上丛可平凡化，情况不同）。此猜想是 猜想 10.1.2 与 (10.1-3) 的前提。]
现在我们处于如下境地：我们期待存在一个范畴 $\cC$，纤维化于“真正的”局部 Langlands 参数空间 $\Loc_{{}^L G}(D^\times)$ 之上，并带有圈群 $G((t))$ 的纤维保持作用。对应于各 $\sigma\in\Loc_{{}^L G}(D^\times)$ 的纤维范畴 $\cC_\sigma$ 应当满足各种尚未明确指出的性质。此外，我们期待 $\cC$ 是带 $G((t))$ 作用的万有范畴。换言之，我们期待 $\Loc_{{}^L G}(D^\times)$ 是 $G((t))$ 的范畴表示的万有参数空间（目前我们还不能更精确地表述这一性质）。圈群局部 Langlands 对应的最终形式大致应当是如下陈述：
\begin{equation}
    \mbox{\fbox{\parbox[c]{4.1cm}{\centering 纤维化于 $\Loc_{{}^L G}(D^\times)$\\ 上的范畴}}}
    \quad\Longleftrightarrow\quad
    \mbox{\fbox{\parbox[c]{3.9cm}{\centering 带 $G((t))$ 作用\\ 的范畴}}}
    \tag{10.1-2}
\end{equation}
其想法是：给定左边的范畴 $\cA$，取 $\cA$ 与万有范畴 $\cC$ 在 $\Loc_{{}^L G}(D^\times)$ 上的“纤维积”，就得到右边的范畴。

% [Tutor Note: p.310 含糊处。(10.1-2) 是纲领性的“最终形式”，不是定理：作者明确说“at the moment we cannot formulate this property more precisely”。所谓“纤维积”是指范畴的基变换（即 (10.1-3) 的笛卡尔性），而不是集合层面的纤维积。阅读时不要把 (10.1-2) 当成已证明的等价。]
现在，我们已经构造出范畴 $\hat{\fg}_{\kappa_c}\text{-}\mathrm{mod}$，它纤维化于 $\Loc_{{}^L G}(D^\times)$ 的一个近亲——即 ${}^L G$-opers 的空间 $\Op_{{}^L G}(D^\times)$——之上，并且带有圈群 $G((t))$ 的纤维保持作用。

$\hat{\fg}_{\kappa_c}\text{-}\mathrm{mod}$ 与猜想中的万有范畴 $\cC$ 之间应当是什么关系？

[Frenkel and Gaitsgory 2006c] 的想法是：第二个纤维化是第一个的“基变换”，即存在笛卡尔方块
\begin{equation}
\begin{array}{ccc}
    \hat{\fg}_{\kappa_c}\text{-}\mathrm{mod} & \longrightarrow & \cC\\[3pt]
    \downarrow & & \downarrow\\[3pt]
    \Op_{{}^L G}(D^\times) & \xrightarrow{\ \alpha\ } & \Loc_{{}^L G}(D^\times)
\end{array}
    \tag{10.1-3}
\end{equation}
它沿两个竖直映射的纤维与 $G((t))$ 的作用交换。换言之，
\[
    \hat{\fg}_{\kappa_c}\text{-}\mathrm{mod}
    \simeq\cC\times_{\Loc_{{}^L G}(D^\times)}\Op_{{}^L G}(D^\times).
\]
于是，$\hat{\fg}_{\kappa_c}\text{-}\mathrm{mod}$ 应当作为对应 (10.1-2) 右端的范畴出现，而与之相配的左端范畴是 $\Op_{{}^L G}(D^\times)$ 上拟凝聚层的范畴。后者经映射 $\Op_{{}^L G}(D^\times)\to\Loc_{{}^L G}(D^\times)$ 纤维化于 $\Loc_{{}^L G}(D^\times)$ 之上。

% [Tutor Note: p.310--311 易错处。(10.1-3) 的“笛卡尔”是对范畴而言：两个基变换范畴等价，且等价与 $G((t))$ 作用相容。这里没有通常意义上的纤维积的严格构造，$\cC$ 也尚无定义，所以 (10.1-3) 只能作为指导原则使用（作者原文亦称 guiding principle）。符号上左端是 $\hat{\fg}_{\kappa_c}\text{-}\mathrm{mod}$、右端是万有范畴 $\cC$，不要与 (10.1-2) 的左右混淆。]
目前我们还没有 $\cC$ 的定义，因此无法把同构写精确。但我们会把它用作指导原则。下面讨论这个猜想的各种推论以及使我们相信它成立的各种证据。

特别地，固定一个在映射 $\alpha$ 的像中的 Langlands 参数 $\sigma\in\Loc_{{}^L G}(D^\times)$（按 猜想 \ref{conj:u14-alpha}，所有 Langlands 参数都在像中）。设 $\chi$ 是 $\sigma$ 的原像 $\alpha^{-1}(\sigma)$ 中的一个 ${}^L G$-oper。那么按上述猜想，范畴 $\hat{\fg}_{\kappa_c}\text{-}\mathrm{mod}_\chi$ 等价于依附于 $\sigma$ 的“未定”Langlands 范畴 $\cC_\sigma$。因此可以把 $\hat{\fg}_{\kappa_c}\text{-}\mathrm{mod}_\chi$ 作为 $\cC_\sigma$ 的定义。\transnote{“未定”（would-be）指 $\cC_\sigma$ 本应是 $\sigma$ 对应的范畴，但此时还没有定义；[Frenkel and Gaitsgory 2006c] 的办法是用 $\hat\fg_{\kappa_c}\text{-}\mathrm{mod}_\chi$（$\chi$ 取 $\alpha^{-1}(\sigma)$ 中任一点）反过来定义它。这个定义只有在断言不同 $\chi$ 给出的范畴典范等价之后才合法，而这正是猜想 10.1.2；在猜想之外，$\cC_\sigma$ 仍只是形式记号。[来源：原书 p.311；[Frenkel and Gaitsgory 2006c]。]}

当然，需要注意：必须保证这个定义不依赖于 $\alpha^{-1}(\sigma)$ 中 $\chi$ 的选取。这意味着对 $\sigma$ 的原像中任意两个 ${}^L G$-opers $\chi$ 与 $\chi'$，相应的范畴 $\hat{\fg}_{\kappa_c}\text{-}\mathrm{mod}_\chi$ 与 $\hat{\fg}_{\kappa_c}\text{-}\mathrm{mod}_{\chi'}$ 应当彼此等价，并且这个等价应当与圈群 $G((t))$ 的作用交换。此外，我们还应期待，当沿纤维 $\alpha^{-1}(\sigma)$ 移动时，这些等价彼此相容。这里我们不去把这一条件写得更精确（但下文 猜想 10.3.10 会解释它对正则 opers 意味着什么）。

即使把相容性问题放在一边，我们也得到如下相当不平凡的猜想。

\begin{conjecture}\label{conj:u14-fiber}
    设 $\chi,\chi'\in\Op_{{}^L G}(D^\times)$ 满足 $\alpha(\chi)=\alpha(\chi')$，即 ${}^L G$-opers $\chi$ 与 $\chi'$ 底层的 $D^\times$ 上平坦 ${}^L G$-丛彼此同构。则范畴 $\hat{\fg}_{\kappa_c}\text{-}\mathrm{mod}_\chi$ 与 $\hat{\fg}_{\kappa_c}\text{-}\mathrm{mod}_{\chi'}$ 之间存在与两个范畴上 $G((t))$ 作用交换的等价。
\end{conjecture}

% [Tutor Note: p.311 背景。猜想 10.1.2 是本章第一个“意外对称性”：同一局部系统的不同 oper 结构（相差的是 Borel 约化）给出等价的临界水平模范畴。它把 opers 空间上的纤维化“下降”到局部系统空间，从而与 (10.1-3) 一致。注意猜想只断言范畴等价并且与 $G((t))$ 作用交换，并未断言自然的函子。]
于是，在追寻局部 Langlands 对应的过程中，我们在临界水平的 $\hat{\fg}$-模范畴 $\hat{\fg}_{\kappa_c}\text{-}\mathrm{mod}$ 的结构中发现了一个意外的对称性。\transnote{猜想 10.1.2 只断言：同一局部系统 $\sigma$ 的不同 oper 结构 $\chi,\chi'$ 给出的两个范畴等价，且等价与 $G((t))$ 作用交换。它不断言等价的自然性，也不断言当 $\chi$ 沿纤维 $\alpha^{-1}(\sigma)$ 变动时这些等价彼此相容（后者是更强的要求，见前段与 10.3 节猜想 10.3.10）；若 (10.1-3) 的笛卡尔方块成立，则 10.1.2 是它的形式推论。经典局部 Langlands 对应的精确陈述（未分歧情形的 Satake 参数、L-包、Frobenius 半单性等）见第 1 章与 10.3 节，本章给出的是其范畴版本。[来源：原书 p.311；第 1 章、10.3 节。]}

\section{Harish--Chandra 范畴}

\begin{keypoints}
    \item \textbf{核心思想：} 把 $G(F)$ 的不可约光滑表示“范畴化”：参数 $\sigma$ 对应范畴 $\cC_\sigma$，其 Grothendieck 群应当“看起来像” $G(F)$ 的不可约光滑表示。
    \item \textbf{数学工具：} 容许性与 $K$-不变向量空间 $R^{\pi(K)}$；Hecke 代数 $H(G(F),K)$ 及其卷积；Harish--Chandra 对 $(\hat{\fg}_{\kappa_c},K)$；范畴 Hecke 代数 $\cH(G((t)),K)$（$K$-等变 D-模）。
    \item \textbf{关键结论：} 弱 $K$-等变与强 $K$-等变之别；(10.2-1)、(10.2-2) 的范畴类比是 D-模的卷积，一般须在导出范畴中进行；仅当 $K=K_0$ 时卷积函子正合，可停留在阿贝尔范畴。
    \item \textbf{记号提醒：} $K_0=G[[t]]$，$K_N$ 为第 $N$ 个同余子群，$I$ 为 Iwahori 子群；$\hat{\fg}_{\kappa_c}\text{-}\mathrm{mod}_\chi^K$ 是 $(\hat{\fg}_{\kappa_c},K)$ Harish--Chandra 模的子范畴，对应 $R^{\pi(K)}$。
\end{keypoints}

% [Tutor Note: p.311 背景。本节是全章的技术核心：把“表示 + K-不变向量 + Hecke 代数作用”这条经典链条逐项范畴化。请记住三组对应：表示 ↔ 范畴（$\cC_\sigma$）；向量 ↔ 对象；$R^{\pi(K)}$ ↔ Harish--Chandra 子范畴。这张对照表在 p.315 给出。]
如第 1 章所解释，圈群 $G((t))$ 的局部 Langlands 对应，应当看作群 $G(F)$ 的局部 Langlands 对应的范畴化，其中 $F$ 是局部非阿基米德域。这意味着，我们试图配给 Langlands 参数 $\sigma\in\Loc_{{}^L G}(D^\times)$ 的、带 $G((t))$ 作用的范畴 $\cC_\sigma$，应当看作“1.1.5 节所讨论的相应局部 Langlands 参数所对应的、$G(F)$ 在复向量空间上的光滑表示”的范畴化。这里我们用“范畴化”一词，是指我们期待范畴 $\cC_\sigma$ 的 Grothendieck 群“看起来像” $G(F)$ 的不可约光滑表示。

% [Tutor Note: p.312 含糊处。“范畴化”在这里只是指导性的期待：$\cC_\sigma$ 的 Grothendieck 群应当与某个不可约光滑表示“相似”，并不是严格意义上的范畴化（没有说 Grothendieck 群等于该表示）。真正被证明的是无分歧情形（10.3 节）里 Harish--Chandra 范畴等价于向量空间范畴。]
本章的目标是把范畴 $\cC_\sigma$ 描述为 $\hat{\fg}_{\kappa_c}$-模范畴。我们先更仔细地考察 $G(F)$ 的不可约光滑表示的结构。

\subsection{\texorpdfstring{$K$}{K}-不变向量的空间}

已知 $G(F)$ 的不可约光滑表示 $(R,\pi)$ 自动是容许的：对任意开紧子群 $K$，例如 1.1.2 节所定义的第 $N$ 个同余子群 $K_N$，$R$ 中 $K$-不变向量的空间 $R^{\pi(K)}$ 有限维。因此，$G(F)$ 的多数不可约光滑表示 $(R,\pi)$ 虽然是无穷维的，却被有限维的 $K$-不变向量子空间 $R^{\pi(K)}$ 滤过，其中 $K$ 取越来越小的开紧子群。空间 $R^{\pi(K)}$ 不带 $G(F)$ 的作用，但带 Hecke 代数 $H(G(F),K)$ 的作用。\transnote{容许性（admissibility）是 $p$-进表示论的基本有限性条件：每个开紧子群 $K$ 的不动向量空间有限维；光滑性则给出 $R=\bigcup_K R^{\pi(K)}$，即 $R$ 是这些有限维子空间的并。注意 $R^{\pi(K)}$ 本身不再是 $G(F)$-模，只带 Hecke 代数 $H(G(F),K)$ 的作用——这正是 §10.2 要把“表示”换成“带 $G((t))$ 作用的范畴”的经典原型。[来源：原书 p.312；Borel--Wallach 2000。]}

按定义，$H(G(F),K)$ 是 $G(F)$ 上紧支集的 $K$ 双不变函数空间；它的代数结构由卷积积
\begin{equation}
    (f_1\star f_2)(g)=\int_{G(F)}f_1(gh^{-1})f_2(h)\,dh
    \tag{10.2-1}
\end{equation}
给出，其中 $dh$ 是 $G(F)$ 上的 Haar 测度，其归一化使子群 $K_0=G(\cO)$ 的体积为 $1$（这里 $\cO$ 是 $F$ 的整数环；例如对 $F=\F_q((t))$ 有 $\cO=\F_q[[t]]$）。代数 $H(G(F),K)$ 通过公式
\begin{equation}
    f\star v=\int_{G(F)}f_1(gh^{-1})(\pi(h)\cdot v)\,dh,
    \qquad v\in R^{\pi(K)},
    \tag{10.2-2}
\end{equation}
作用在空间 $R^{\pi(K)}$ 上。

% [Tutor Note: p.312 易错处 + 疑似笔误。(10.2-2) 左端写 $f\star v$，右端被积函数却印作 $f_1(gh^{-1})$；按 (10.2-1) 的记号，这里应为 $f(gh^{-1})$（没有下标 1）。译文照原书保留，报告已登记。另注意 Haar 测度的归一化是 $\mathrm{vol}(K_0)=1$，这决定了 Hecke 代数单位元是 $K_0$ 的特征函数，换归一化会改变 (10.2-2) 的系数。]
研究 $K$-不变向量的空间及其 $H(G(F),K)$-模结构，是分析群 $G(F)$（$F=\F_q((t))$）的表示的有效工具。

在圈群的范畴局部 Langlands 对应中，能否找到类似的结构？

\subsection{等变模}

% [Tutor Note: p.312--313。本节把“K-不变向量”翻译成“范畴的 K-等变性”：对象 $(M,\rho)$ 在 $K$ 下稳定，等价于给出一族与 $G((t))$ 作用交换的同构 $T_k^M$。注意 $K$ 在 $G((t))$ 中并不紧（对通常拓扑），“紧子群”只是类比说法；真正需要的是 $\mathfrak{k}$ 的作用可积。]
在范畴情形，群 $G(F)$ 的表示 $(R,\pi)$ 被一个带 $G((t))$ 作用的范畴替代，例如 $\hat{\fg}_{\kappa_c}\text{-}\mathrm{mod}_\chi$。$G(F)$ 的开紧子群对圈群 $G((t))$ 有明显的类比物（当然，它们对 $G((t))$ 上的通常拓扑并不紧）。例如有“极大紧子群”$K_0=G[[t]]$，或更一般地，第 $N$ 个同余子群 $K_N$，其元素模 $t^N\C[[t]]$ 同余于 $1$。另一个重要例子是 Iwahori 子群的类比物。它是 $G[[t]]$ 的子群，记作 $I$，其元素 $g(t)$ 满足：$g(t)$ 在 $0$ 处的值 $g(0)$ 属于固定的 Borel 子群 $B\subset G$。

现在，在范畴情形，$G(F)$ 的表示中的向量的类比物是范畴的对象，即光滑 $\hat{\fg}_{\kappa_c}$-模 $(M,\rho)$，其中 $\rho:\hat{\fg}_{\kappa_c}\to\End M$。因此对上述类型的子群 $K\subset G((t))$，$G(F)$ 的表示中 $K$-不变向量的类比物，是在 $K$ 的作用下稳定的光滑 $\hat{\fg}_{\kappa_c}$-模 $(M,\rho)$。回忆 1.3.6 节：对任意 $g\in G((t))$ 有新的 $\hat{\fg}_{\kappa_c}$-模 $(M,\rho_g)$，其中 $\rho_g(x)=\rho(\operatorname{Ad}_g(x))$。称 $(M,\rho)$ 在 $K$ 下稳定，或称 $(M,\rho)$ 是弱 $K$-等变的，若对一切 $k\in K$，$(M,\rho)$ 与 $(M,\rho_k)$ 之间存在相容的同构系统。更精确地说，这意味着对每个 $k\in K$ 存在线性映射 $T_k^M:M\to M$，使得
\[
    T_k^M\rho(x)(T_k^M)^{-1}=\rho(\operatorname{Ad}_k(x))
\]
对一切 $x\in\hat{\fg}_{\kappa_c}$ 成立，并且
\[
    T_1^M=\mathrm{Id}_M,\qquad T_{k_1}^M T_{k_2}^M=T_{k_1k_2}^M.
\]
于是 $M$ 成为群 $K$ 的表示。\footnote{一般地，把最后一个条件修改为允许非平凡的 $2$-上循环、从而允许 $K$ 的非平凡中心扩张是合理的；但在我们关心的情形，$K$ 没有非平凡的中心扩张。} 考虑李代数 $\mathfrak{k}=\Lie K$ 在 $M$ 上相应的表示。设嵌入 $\mathfrak{k}\hookrightarrow\fg((t))$ 提升为 $\mathfrak{k}\hookrightarrow\hat{\fg}_{\kappa_c}$（即中心扩张上循环在 $\mathfrak{k}$ 上平凡）。例如，对含于 $K_0=G[[t]]$ 的任意子群或其共轭，这都是成立的。于是我们还有由 $\rho$ 限制得到的 $\mathfrak{k}$ 在 $M$ 上的表示。一般地，两个表示不必相同。若它们相同，则称模 $M$ 是强 $K$-等变的，或简称 $K$-等变的。

% [Tutor Note: p.313 易错处。弱 $K$-等变只要求每个 $k$ 给出抽象的线性同构 $T_k^M$；强 $K$-等变还要求这个同构与李代数 $\mathfrak{k}$ 通过 $\rho$ 已给出的作用一致。判据是：$T_k^M$ 与 $\rho(\mathfrak{k})$ 的作用是否吻合。另注意 $T_{k_1}^M T_{k_2}^M=T_{k_1k_2}^M$ 的乘法顺序按原书（若把 $k\mapsto T_k$ 看成反表示，则顺序会相反）；这里按原文照录，核对时不要“顺手”交换。]
对 $(\hat{\fg}_{\kappa_c},K)$ 本身是 Harish--Chandra 对的一个例子。所谓 Harish--Chandra 对，是指由李代数 $\fg$ 与李群 $H$ 组成的对 $(\fg,H)$，其中 $H$ 的李代数含于 $\fg$。因此 $K$-等变的 $\hat{\fg}_{\kappa_c}$-模称为 $(\hat{\fg}_{\kappa_c},K)$ Harish--Chandra 模。它们是（光滑的）$\hat{\fg}_{\kappa_c}$-模，且李代数 $\Lie K\subset\hat{\fg}_{\kappa_c}$ 的作用可指数化为 $K$ 的作用（我们假定 $K$ 连通）。我们用 $\hat{\fg}_{\kappa_c}\text{-}\mathrm{mod}^K$ 与 $\hat{\fg}_{\kappa_c}\text{-}\mathrm{mod}_\chi^K$ 分别记 $\hat{\fg}_{\kappa_c}\text{-}\mathrm{mod}$ 与 $\hat{\fg}_{\kappa_c}\text{-}\mathrm{mod}_\chi$ 中对象为 $(\hat{\fg}_{\kappa_c},K)$ Harish--Chandra 模的满子范畴。\transnote{Harish--Chandra 对 $(\fg,H)$ 的一般定义还要求 $H$ 通过伴随作用作用于 $\fg$，且 $\operatorname{Lie}H$ 的作用是它的微分；这里 $H=K\subset G((t))$ 的伴随作用给出这一结构。(g,K)-模的等价说法是：$M$ 是 $\fg$-模，$\operatorname{Lie}K$ 的作用可指数化为 $K$ 的作用，并满足 $T_k^M\rho(x)(T_k^M)^{-1}=\rho(\operatorname{Ad}_k x)$——这正是“强 $K$-等变”。要求 $K$ 连通，是为了让 $\operatorname{Lie}K$ 的作用总能指数化。[来源：原书 p.313；Borel--Wallach 2000。]}

我们规定：范畴 $\hat{\fg}_{\kappa_c}\text{-}\mathrm{mod}_\chi$ 中 $K$-不变向量的类比物是 $(\hat{\fg}_{\kappa_c},K)$ Harish--Chandra 模。因此，范畴 $\hat{\fg}_{\kappa_c}\text{-}\mathrm{mod}_\chi$ 应当看作群 $G(F)$ 的光滑不可约表示 $(R,\pi)$ 的类比物，而范畴 $\hat{\fg}_{\kappa_c}\text{-}\mathrm{mod}_\chi^K$ 是 $K$-不变向量空间 $R^{\pi(K)}$ 的类比物。

\subsection{范畴 Hecke 代数}

接着讨论 Hecke 代数 $H(G(F),K)$ 的范畴类比物。

回忆 $H(G(F),K)$ 是 $G(F)$ 上紧支集的 $K$ 双不变函数代数。我们把它实现为 $G(F)/K$ 上左 $K$-不变的紧支集函数代数。在 1.3.3 节我们已讨论过齐性空间（如 $G(F)/K$）上函数代数的范畴化问题。我们的结论是：当 $G(F)$ 换成复圈群 $G((t))$ 时，这个代数的范畴类比物是 $G((t))/K$ 上 D-模的范畴。更精确地说，这个商具有 ind-概形结构，即沿闭嵌入的有限维代数簇的正向极限。这样的 ind-概形上（右）D-模的适当概念在 [Beilinson and Drinfeld 1997] 中给出（另见 [Frenkel and Gaitsgory 2004; Frenkel and Gaitsgory 2006c]）。作为 $G(F)/K$ 上左 $K$-不变函数代数的范畴类比物，我们取 ind-概形 $G((t))/K$ 上 $K$-等变 D-模的范畴 $\cH(G((t)),K)$（关于 $K$ 在 $G((t))/K$ 上的左作用）。我们称它为与 $K$ 相联系的范畴 Hecke 代数。\transnote{$G((t))/K$ 只是 ind-概形（有限维闭子概形的正向极限），通常的 D-模定义不能直接照搬：[Beilinson and Drinfeld 1997] 采取“在所有有限维子概形上相容地给出一族 D-模”的方式来定义。$\cH(G((t)),K)$ 的乘法是 D-模的卷积（由群乘法诱导的 $*$-或 $!$-推前），这正是下文要讨论推前选择与导出范畴的原因。[来源：原书 p.314；[Beilinson and Drinfeld 1997]；[Frenkel and Gaitsgory 2006c]。]}

模仿公式 (10.2-1)，容易定义范畴 $\cH(G((t)),K)$ 中两个对象的卷积。即，把这个公式解释为“拉回与函数积分”两种运算的复合。然后把同样的运算施加于 D-模，把积分理解为推前。然而这里会遇到两个问题。第一个问题是：对一般群 $K$，涉及的态射不是紧合的，因此必须在 $*$-推前与 $!$-推前之间作选择。不过，若 $K$ 满足 $I\subset K\subset G[[t]]$，这个问题不会出现；这正是我们主要关心的情形。第二个、也更严重的问题是：一般地推前不是正合函子，因此两个 D-模的卷积不再是 D-模，而是一个复形，更精确地说，是 $G((t))/K$ 上 D-模的相应 $K$-等变（有界）导出范畴 $D^b(G((t))/K)^K$ 的对象。这里我们不写出这个范畴的确切定义，有兴趣的读者可参看 [Beilinson and Drinfeld 1997] 与 [Frenkel and Gaitsgory 2006c]。例外是子群 $K_0=G[[t]]$ 的情形：此时卷积函子正合，因此可以只考虑 $G((t))/K_0$ 上 $K_0$-等变 D-模的阿贝尔范畴。

% [Tutor Note: p.314 易错处。这里有两个技术要点：其一，$*$-推前与 $!$-推前在 proper 态射时一致，$I\subset K\subset G[[t]]$ 保证了所需的 properness（见 [Beilinson and Drinfeld 1997] 的约定）；其二，卷积的“正合性”依赖于 $K$，一般必须在导出范畴 $D^b(G((t))/K)^K$ 中工作。作者特别指出只有 $K_0$ 情形可停留在阿贝尔范畴，这与 10.3 节无分歧情形的计算直接相关。]
范畴 $D^b(G((t))/K)^K$ 有幺半结构，因此作用在 $(\hat{\fg}_{\kappa_c},K)$ Harish--Chandra 模的导出范畴上（确切定义同样请参看 [Beilinson and Drinfeld 1997; Frenkel and Gaitsgory 2006c]）。在 $K=K_0$ 的特殊情形，我们可以只考虑相应的阿贝尔范畴。这个作用应当看作 $H(G(F),K)$ 在 $K$-不变向量空间 $R^{\pi(K)}$ 上作用的范畴类比物。

我们的最终目标是理解与“局部 Langlands 参数” $\sigma\in\Loc_{{}^L G}(D^\times)$ 相联系的“局部 Langlands 范畴” $\cC_\sigma$。现在对范畴 $\cC_\sigma$ 我们有了候选：即范畴 $\hat{\fg}_{\kappa_c}\text{-}\mathrm{mod}_\chi$，其中 $\sigma=\alpha(\chi)$。因此 $\hat{\fg}_{\kappa_c}\text{-}\mathrm{mod}_\chi$ 应当看作 $G(F)$ 的光滑表示 $(R,\pi)$ 的范畴化。相应地，$(\hat{\fg}_{\kappa_c},K)$ Harish--Chandra 模的范畴 $\hat{\fg}_{\kappa_c}\text{-}\mathrm{mod}_\chi^K$ 应当看作 $R^{\pi(K)}$ 的范畴化。范畴 Hecke 代数 $\cH(G((t)),K)$ 作用在这个范畴上。实际上，这个作用并不保持阿贝尔范畴 $\hat{\fg}_{\kappa_c}\text{-}\mathrm{mod}_\chi^K$；$\cH(G((t)),K)$ 作用在相应的导出范畴上。\footnote{这是局部 Langlands 对应应当把局部系 $\sigma\in\Loc_{{}^L G}(D^\times)$ 配上一个三角范畴（并带 $G((t))$ 作用），而不是阿贝尔范畴的第一个迹象。下文所考察的例子将显示，这些三角范畴带有不同的 $t$-结构。换言之，它们可以解释为不同阿贝尔范畴的导出范畴，而我们期待其中之一就是 $\hat{\fg}_{\kappa_c}\text{-}\mathrm{mod}_\chi$。}\transnote{三角范畴上的 $t$-结构（[Beilinson, Bernstein and Deligne 1982]）给出“阿贝尔范畴 + 导出范畴”的相对概念：同一个三角范畴可以带不同的 $t$-结构，其心是不同的阿贝尔范畴。脚注 2 的意思是：局部 Langlands 对应的目标应是带 $G((t))$ 作用的三角范畴，$\hat\fg_{\kappa_c}\text{-}\mathrm{mod}_\chi$ 只是它在某个 $t$-结构下的心；几何 Langlands 中 Hecke 本征层取 perverse $t$-结构，正是同一现象的另一例。[来源：原书 p.315 脚注 2；[Beilinson, Bernstein and Deligne 1982]。]}

% [Tutor Note: p.315 背景。脚注 2 是本章的一个重要提示：局部 Langlands 对应的目标范畴本质上是三角范畴，阿贝尔范畴（如 $\hat{\fg}_{\kappa_c}\text{-}\mathrm{mod}_\chi$）只是它的一种 $t$-结构选择。这与几何 Langlands 对应中“Hecke 本征层是反常层（perverse sheaf）”的现象一脉相承：perverse $t$-结构与“模范畴”的 $t$-结构往往是两个不同的选择。]
我们把这一类比总结成下表。

\begin{center}
\begin{tabular}{p{0.42\linewidth}p{0.42\linewidth}}
\toprule
经典理论 & 几何理论 \\
\midrule
$G(F)$ 在向量空间 $R$ 上的表示 & $G((t))$ 在范畴 $\hat{\fg}_{\kappa_c}\text{-}\mathrm{mod}_\chi$ 上的表示 \\
$R$ 中的向量 & $\hat{\fg}_{\kappa_c}\text{-}\mathrm{mod}_\chi$ 的对象 \\
$R$ 的 $K$-不变向量子空间 $R^{\pi(K)}$ & $(\hat{\fg}_{\kappa_c},K)$ Harish--Chandra 模的子范畴 $\hat{\fg}_{\kappa_c}\text{-}\mathrm{mod}_\chi^K$ \\
Hecke 代数 $H(G(F),K)$ 作用在 $R^{\pi(K)}$ 上 & 范畴 Hecke 代数 $\cH(G((t)),K)$ 作用在 $\hat{\fg}_{\kappa_c}\text{-}\mathrm{mod}_\chi^K$ 上 \\
\bottomrule
\end{tabular}
\end{center}

% [Tutor Note: p.315。这张表是全章的地图：左列是 1.1--1.3 节的经典对象，右列是本章的范畴类比。注意最后一行的“作用”在右边必须理解为导出范畴上的作用（见脚注 2），而表中未标出这一点，引用时请留意。]
现在可以通过研究 Harish--Chandra 模范畴 $\hat{\fg}_{\kappa_c}\text{-}\mathrm{mod}_\chi^K$、并把它们的结构与已知情形中 $G(F)$ 的光滑表示的 $K$-不变向量空间 $R^{\pi(K)}$ 的结构相比较，来检验我们关于局部 Langlands 对应的提法。另一种可能是把 猜想 \ref{conj:u14-fiber} 应用于 Harish--Chandra 模范畴来检验它。

在下一节我们考虑“极大紧子群”$K_0=G[[t]]$ 的情形。我们将证明范畴 $\hat{\fg}_{\kappa_c}\text{-}\mathrm{mod}_\chi^{K_0}$ 的结构与关于 $G(F)$ 无分歧表示的经典结果相容。然后我们在 10.4 节处理更复杂的 Iwahori 子群 $I$ 的情形。在那里我们也会看到，[Frenkel and Gaitsgory 2006c] 的猜想与结果同关于 $G(F)$ 的带 Iwahori 不动向量表示的已知结果是一致的。
